%!TEX root = ../main.tex

% ==================================================================
\section{Lecture~2: Periods and GV invariants from toric data}\label{sec:L2}

\medskip

\noindent\textit{Anchor:} TASI \cite{McAllister:2025qwq},
\S\S\,3--6.

\medskip



\noindent\textit{Guiding question.}
Lecture~1 ended with a flat moduli space: the Calabi--Yau metric
can move in continuous families, and at leading order the associated
four-dimensional scalar fields have $V=0$.  Lecture~2 asks how one
computes the first ingredient that turns this geometric moduli space
into an explicit function in the four-dimensional theory.  The answer
is a period calculation.  Quantised three-form fluxes choose an
integral linear combination of periods of the holomorphic three-form,
so that
\begin{equation}
    W_{\rm flux}
    =
    \int_X \Omega\wedge G_3\, .
\end{equation}
The target of the lecture is therefore deliberately narrow: compute
$W_{\rm flux}$ by computing periods.  Mirror symmetry converts the
period problem near a large complex structure point of $X$ into a
large volume calculation on the mirror $\widetilde X$.  For toric
hypersurfaces, the relevant data are finite and combinatorial: the
reflexive polytope, the GLSM charge matrix, the intersection numbers, and
the genus-zero Gopakumar--Vafa invariants \cite{Gopakumar:1998ii,Gopakumar:1998jq} that organise the
worldsheet instanton contributions to the mirror prepotential \cite{Candelas:1993dm,Candelas:1994hw}.  The final
section then uses this period machinery in one application, namely the
perturbatively flat vacuum construction of \cite{Demirtas:2019sip}.

The effective field theory and compactification background follow the
conventions of \cite{McAllister:2025qwq}, especially the material on
flux compactifications and periods.  The expanded discussion of period
computation and genus-zero GV data below follows the computational
mirror-symmetry framework of
\cite{Demirtas:2023als,RiosTascon:2023thesis}.  Notice the order of
logic: we first compute periods and the flux superpotential; only at
the end do we discuss a particular flux-vacuum application.

\medskip
\noindent\textit{Historical and computational thread.}
The bridge from Calabi--Yau geometry to numbers in the
four-dimensional theory is mirror symmetry. The quintic calculation
of Candelas, de~la~Ossa, Green and Parkes
\cite{Candelas:1990rm}, Witten's topological field theory
formulation \cite{Witten:1991zz}, the two-parameter examples
\cite{Candelas:1993dm,Candelas:1994hw}, and the
Hosono--Klemm--Theisen--Yau (HKTY) period/mirror-map construction
\cite{Hosono:1993qy,HKTY2} provide the historical route from
periods on the mirror to instanton-corrected prepotentials.
The SYZ and Hori--Vafa perspectives give complementary geometric
and physical explanations of mirror symmetry
\cite{Strominger:1996it,Hori:2000kt,Hori:2003ic}, although the
computations below use the period/GKZ side of the story.
Batyrev's toric mirror construction \cite{Batyrev:1993oya} and the
large complex structure compactification picture
\cite{morrison1993compactifications} explain why this is the
natural language for hypersurfaces in toric varieties. The
Gopakumar--Vafa invariants \cite{Gopakumar:1998ii,Gopakumar:1998jq}
then give the integral curve-counting data that organise the
genus-zero instanton expansion used below. Computational mirror
symmetry in the sense of \cite{Demirtas:2023als} is the
high-dimensional continuation of this story: it asks how to compute
the period vector, mirror map and GV data reliably when the Picard
rank is large and the Mori cone need not be simplicial.





\subsection{What has to be computed?}\label{ssec:L2-computational-targets}

The first point of the lecture is not yet mirror symmetry.  It is to
identify the four-dimensional quantity for which the geometry of $X$
will become useful.  A compactification turns continuous families of
Calabi--Yau metrics and background fields into scalar fields in four
spacetime dimensions.  Locally on field space we write these fields as
\begin{equation}\label{eq:L2-field-content}
    \Phi^I=(\tau,z^a,T_i)\, .
\end{equation}
Here $z^a$, with $a=1,\ldots,h^{2,1}(X)$, are local coordinates on
complex structure moduli space; $T_i$, with
$i=1,\ldots,h^{1,1}(X)$, are complexified K\"ahler moduli; and
$\tau$ is the axio-dilaton.  The ten-dimensional origin of $\tau$ and of the
three-form fluxes is recalled in \S\ref{ssec:L2-IIB}; for the present
section only the four-dimensional role of these fields matters.

For slowly varying fields, the scalar part of the four-dimensional
Einstein-frame action has the schematic form
\begin{equation}\label{eq:L2-schematic-4d-action}
    S_4
    =
    \int \mathrm{d}^4x\,\sqrt{-g_4}\,
    \left(
        \frac12 R_4
        -K_{I\bar J}(\Phi,\bar\Phi)\,
          \partial_\mu\Phi^I\partial^\mu\bar\Phi^{\bar J}
        -V(\Phi,\bar\Phi)
        +\cdots
    \right)\, .
\end{equation}
This is the lecture's basic organising equation.  The coefficient
$K_{I\bar J}$ is the metric on the scalar field space, so it tells us
what ``gradient'' means in the equations of motion.  The scalar
potential $V$ is the energy density for constant fields, and its
critical points are the candidate four-dimensional vacua.
More explicitly, the field equations have the schematic form
\begin{equation}\label{eq:L2-scalar-EOM-schematic}
    \Box\Phi^I
    +\Gamma^{I}_{JK}\,\partial_\mu\Phi^J\partial^\mu\Phi^K
    =
    K^{I\bar J}\,\partial_{\bar J}V\, ,
\end{equation}
up to the conventional sign determined by the spacetime signature.
For constant scalar fields this reduces to the critical-point equations
\begin{equation}\label{eq:L2-critical-point-equations}
    \partial_I V=0\, .
\end{equation}
Thus the vacuum problem asks for two pieces of data: the field-space
metric and the scalar potential.

The field-space metric is the first place where the geometry of the
Calabi--Yau enters.  In the supersymmetric compactifications used in
these notes, $K_{I\bar J}$ is a K\"ahler metric,
\begin{equation}\label{eq:L2-Kahler-metric}
    K_{I\bar J}
    =
    \partial_I\partial_{\bar J}K\, ,
\end{equation}
for a real function $K$, the K\"ahler potential.  On the complex
structure moduli space this metric is the Weil--Petersson metric.
The Hodge-theoretic formula is
\begin{equation}\label{eq:L2-Kcs-intro}
    K_{\rm cs}
    =
    -\log\!\Bigl(-\mathrm{i}\int_X\Omega\wedge\bar\Omega\Bigr)\, ,
    \qquad
    G^{\rm cs}_{a\bar b}
    =
    \partial_a\partial_{\bar b}K_{\rm cs}\, .
\end{equation}
This single display contains the reason periods are unavoidable:
$\int_X\Omega\wedge\bar\Omega$ is computed from the periods of the
holomorphic three-form.  Consequently, periods determine not only a
holomorphic superpotential later on, but also the metric and
connection used in the equations of motion.

The K\"ahler-moduli metric is given by
\begin{equation}
 G^{\rm K}_{ij}
    =
    \frac{1}{4\Vol(X,J)}
    \int_X\omega_i\wedge\star_6\omega_j\, ,
\end{equation}
for a basis $\omega_i\in H^{1,1}(X)$, $i=1,\ldots,h^{1,1}(X)$. As
shown in \S\ref{ssec:L2-WP}, it is controlled by the classical volume
function.  If
\begin{equation}
    J=t^i\omega_i
\end{equation}
is the K\"ahler form expanded in this basis, then
\begin{equation}\label{eq:L2-volume-intro}
    \Vol(X,J)
    =
    \frac{1}{3!}\int_X J\wedge J\wedge J
    =
    \frac{1}{6}\kappa_{ijk}t^it^jt^k\, ,
    \qquad
    \kappa_{ijk}=\int_X\omega_i\wedge\omega_j\wedge\omega_k\, .
\end{equation}
The real parts of the complexified K\"ahler coordinates used below are
divisor volumes:
\begin{equation}\label{eq:tauAdef}
    \tau_i
    =
    \mathrm{Re}\,(T_i)
    =
    \frac12\int_X\omega_i\wedge J\wedge J
    =
    \frac{\partial\Vol}{\partial t^i}
    =
    \frac12\kappa_{ijk}t^jt^k\, .
\end{equation}
At leading order in $\alpha'$ and $g_s$, the K\"ahler potential used
for the compactification examples is therefore
\begin{equation}\label{eq:K-full}
    K(\tau,z,T;\bar\tau,\bar z,\bar T)
    =
    -\log\!\bigl(-\mathrm{i}(\tau-\bar\tau)\bigr)
    -\log\!\Bigl(-\mathrm{i}\int_X\Omega\wedge\bar\Omega\Bigr)
    -2\log\Vol(T,\bar T)\, .
\end{equation}
The first term is the axio-dilaton metric, the second is the
Weil--Petersson metric on complex structure moduli space, and the
third is the large-volume K\"ahler-moduli metric.  The more detailed
metric interpretation is given in \S\ref{ssec:L2-WP}; the important
message here is that each term in $K$ is computable from geometric
data.

We can now introduce the scalar potential in the form used later.  In
four-dimensional $\mathcal N=1$ supergravity, the $F$-term potential is
\begin{equation}\label{eq:L2-Fterm-potential}
    V
    =
    \mathrm{e}^K
    \left(K^{I\bar J}D_IW\,\overline{D_JW}-3|W|^2\right)\, ,
    \qquad
    D_IW
    =
    \partial_IW+(\partial_IK)W\, .
\end{equation}
Here $W$ is a holomorphic superpotential.  Thus the compactification
problem becomes concrete only after we can compute $K$, $W$, and the
region in which the formulae are reliable.  Lecture~2 focuses on the
part of this problem that toric mirror symmetry computes most
directly: the flux superpotential
\begin{equation}\label{eq:L2-Wflux-preview}
    W_{\rm flux}(\tau,z)
    =
    \int_X\Omega(z)\wedge G_3(\tau)\, .
\end{equation}
The same periods of $\Omega$ that occur in $K_{\rm cs}$ also occur in
this integral.  Quantised fluxes choose an integral linear combination
of those periods.  Therefore the computational target for the rest of
the lecture is
\begin{equation}\label{eq:L2-computational-target}
    \text{toric data}
    \quad\Longrightarrow\quad
    \text{period vector of }\Omega
    \quad\Longrightarrow\quad
    K_{\rm cs}\text{ and }W_{\rm flux}\, ,
\end{equation}
which we then use to find explicit solutions to
\begin{equation}
D_\tau W_{\rm flux}=D_a W_{\rm flux}=0\, .
\end{equation}
For generic fluxes, the last equations are transcendental and are
solved numerically.  The final section of the lecture explains why
perturbatively flat vacua provide a special class in which the
polynomial large complex structure part can be solved analytically,
and the GV-controlled instanton terms then provide the first
non-trivial lifting.






\subsection{Aside: Weil--Petersson metrics on moduli spaces of Calabi--Yau threefolds\texorpdfstring{$^{\ast}$}{*}}\label{ssec:L2-WP}



The four-dimensional $\mathcal{N}=1$ effective theory associated
with a flux compactification on a Calabi--Yau threefold $X$ inherits
two distinct metrics on its scalar manifold: one on the complex
structure moduli space $\mathcal{M}_{\rm cs}(X)$, the other on the
complexified K\"ahler moduli space $\mathcal{M}_{\rm K}(X)$. Both
are Weil--Petersson metrics in the classical
differential-geometric sense, and both admit two complementary
descriptions:
\begin{itemize}
    \item an $L^2$-pairing of metric variations of the underlying
    Ricci-flat metric on $X$;
    \item a Hodge-theoretic or intersection-theoretic description
    in terms of the holomorphic three-form $\Omega$ on the complex
    structure side, or the K\"ahler form $J$ on the K\"ahler side.
\end{itemize}
The point of the second description is computational: it replaces an
unknown Ricci-flat metric by periods and intersection numbers.  The
unobstructedness of complex structure deformations is the
Tian--Todorov theorem \cite{tian1987smoothness,todorov1989weil};
the equality between the metric-variation pairing and the
Hodge-theoretic formula below is the standard
variation of Hodge structure description of the Weil--Petersson
metric.  

\subsubsection*{Complex structure moduli space \texorpdfstring{$\mathcal{M}_{\rm cs}(X)$}{Mcs}}

Let $g^{(0)}_{m\bar n}$ denote the Ricci-flat K\"ahler metric on
$X$ supplied by Yau's theorem.  Complex structure deformations,
parametrised by local coordinates $z^a$ with
$a=1,\ldots,h^{2,1}(X)$, appear as variations of the metric of pure
type $(2,0)$ and $(0,2)$.  The natural $L^2$ inner product of such
variations is the Weil--Petersson metric
\cite{Candelas:1990pi,Candelas:1989ug}
\begin{equation}\label{eq:WP-cs-L2}
    2\,G^{\rm cs}_{a\bar b}\,
    \mathrm{d}z^a\,\mathrm{d}\bar z^{\bar b}
    =
    \frac{1}{2\mathcal{V}}
    \int_X
    g^{m\bar n}g^{k\bar l}\,
    \delta g_{mk}\,\delta g_{\bar n\bar l}\,
    \sqrt{g}\,\mathrm{d}^6y\, .
\end{equation}
Here $\mathcal{V}=\int_X\sqrt{g}\,\mathrm{d}^6y$ is the total
volume of $X$.  The factor of $2$ on the left records the convention
$\mathrm{d}s^2=2\,G_{a\bar b}\,\mathrm{d}z^a\mathrm{d}\bar z^{\bar b}$;
normalisations differ slightly in the literature.

Via the Calabi--Yau condition, each complex structure deformation
corresponds to a harmonic $(2,1)$-form.  
Concretely, the variation of $\Omega$ satisfies Griffiths
transversality,
\begin{equation}\label{eq:Griffiths}
    \partial_a\Omega
    =
    k_a(z)\,\Omega+\chi_a\, ,
    \qquad
    \chi_a\in H^{2,1}(X)\, .
\end{equation}
Equivalently, after choosing the normalisation determined by
$K_{\rm cs}$ below,
\begin{equation}
    \chi_a
    =
    D_a\Omega
    =
    \partial_a\Omega+(\partial_aK_{\rm cs})\Omega\, .
\end{equation}
The Hodge-theoretic form of the Weil--Petersson metric is
\begin{equation}\label{eq:WP-cs-Hodge}
    G^{\rm cs}_{a\bar b}
    =
    -\,\frac{\displaystyle\int_X
        \chi_a\wedge\bar\chi_{\bar b}}
        {\displaystyle\int_X
        \Omega\wedge\bar\Omega}\, .
\end{equation}
The denominator is the natural Hermitian metric on the Hodge line
bundle whose fibre is $H^{3,0}(X)$.  In particular,
\begin{equation}\label{eq:K-cs-WP}
    K_{\rm cs}
    =
    -\log\!\Bigl(
        -\mathrm{i}\int_X\Omega\wedge\bar\Omega
    \Bigr)\, ,
    \qquad
    G^{\rm cs}_{a\bar b}
    =
    \partial_a\partial_{\bar b}K_{\rm cs}\, .
\end{equation}
A short verification uses \eqref{eq:Griffiths} and
$\int_X\Omega\wedge\chi_a=0$ to obtain
\begin{equation}
    \partial_a\partial_{\bar b}K_{\rm cs}
    =
    -\,\frac{\int_X\chi_a\wedge\bar\chi_{\bar b}}
            {\int_X\Omega\wedge\bar\Omega}\, .
\end{equation}
Thus the period calculation is already a calculation of the
geometry of the equations of motion: the same periods determine the
holomorphic function $W_{\rm flux}$ and the metric and connection
appearing in $D_a W_{\rm flux}$.



\subsubsection*{K\"ahler moduli space \texorpdfstring{$\mathcal{M}_{\rm K}(X)$}{MK}}

The K\"ahler class $[J]$ of the Ricci-flat metric is parametrised,
in a basis $\{\omega_i\}_{i=1}^{h^{1,1}(X)}$ of
$H^{1,1}(X,\mathbb{Z})$, by real moduli $t^i$ lying in the K\"ahler
cone $\KC(X)\subset H^{1,1}(X,\mathbb{R})$:
\begin{equation}
    J=\sum_i t^i\omega_i\in\KC(X)\, .
\end{equation}
The intersection numbers
\begin{equation}
    \kappa_{ijk}
    =
    \int_X\omega_i\wedge\omega_j\wedge\omega_k
    \in\mathbb{Z}
\end{equation}
and the K\"ahler cone determine the classical K\"ahler-side
geometry.  

The $L^2$ inner product of K\"ahler metric variations
$\delta g_{m\bar n}=(\omega_i)_{m\bar n}\,\delta t^i$ induces the
Weil--Petersson metric on K\"ahler moduli space
\cite{Strominger:1990pd,Candelas:1990pi}
\begin{equation}\label{eq:WP-K-L2}
    G^{\rm K}_{ij}
    =
    \frac{1}{4\mathcal{V}}
    \int_X\omega_i\wedge\star_6\omega_j\, .
\end{equation}
The $1/(4\mathcal{V})$ prefactor here is fixed by the
Hodge-star action on $H^{1,1}$ and the volume normalisation.
Applying the Hodge star on $(1,1)$-forms and using the hard-Lefschetz
decomposition gives the closed expression
\begin{equation}\label{eq:WP-K-explicit}
    G^{\rm K}_{ij}
    =
    -\,\frac{\kappa_{ijk}t^k}{4\mathcal{V}}
    +
    \frac{\mathcal{T}_i\mathcal{T}_j}{4\mathcal{V}^2},
    \qquad
    \mathcal{T}_i
    =
    \frac12\int_XJ\wedge J\wedge\omega_i
    =
    \frac12\kappa_{ijk}t^jt^k\, .
\end{equation}
Here $\mathcal{T}_i$ is the volume of the divisor dual to
$\omega_i$; in the main flow it is denoted
$\tau_i=\mathrm{Re}\,T_i$ after choosing the corresponding divisor
basis.  The corresponding tree-level K\"ahler potential is
\begin{equation}\label{eq:K-tree-WP}
    K_{\rm tree}
    =
    -2\log\mathcal{V}
    =
    -2\log\!\Bigl(
        \frac{1}{6}\kappa_{ijk}t^it^jt^k
    \Bigr)\, .
\end{equation}
Differentiating $K_{\rm tree}$ with respect to the complex chiral
coordinates $T_i$ reproduces \eqref{eq:WP-K-explicit}, up to the
Jacobian factors from the change of variables $t^i\mapsto T_i$.


\subsection{Aside: Minimal Type IIB and orientifold dictionary\texorpdfstring{$^{\ast}$}{*}}\label{ssec:L2-IIB}\label{ssec:L2-orientifold}

Only a small part of the ten-dimensional Type IIB
vocabulary is needed for the computations in this lecture.  The
bosonic fields that enter the formulae below are
\begin{equation}\label{eq:axio-dilaton}
    g_{MN},\qquad
    \tau=C_0+\mathrm{i} \mathrm{e}^{-\phi},\qquad
    B_2,\ C_2,\ C_4\, .
\end{equation}
Here $g_{MN}$ is the ten-dimensional metric, $\phi$ is the dilaton,
$C_0$ is the RR axion, $B_2$ and $C_2$ are two-form gauge
potentials, and $C_4$ is a four-form gauge potential.  Their
field strengths are
\begin{equation}
    H_3=\mathrm{d}B_2,\qquad
    F_3=\mathrm{d}C_2,\qquad
    G_3=F_3-\tau H_3\, ,
\end{equation}
with the modified self-dual five-form built from $C_4$ not needed
explicitly in this lecture.  Thus the symbol $\tau$ in the
four-dimensional theory is the zero mode of the axio-dilaton, while
$F_3$ and $H_3$ are the quantised three-form fluxes whose periods
enter $W_{\rm flux}$.

\begin{definition}[O3/O7 orientifold of $X$]\label{def:orientifold}
An \emph{O3/O7 orientifold} of $X$ is the quotient by a
holomorphic involution $\sigma:X\to X$ with $\sigma^2=\mathrm{id}$
and $\sigma^*\Omega=-\Omega$, combined in the physical theory with
worldsheet orientation reversal and fermion parity.  The fixed locus
of $\sigma$ is a disjoint union of isolated points, called
O3-planes, and holomorphic divisors, called O7-planes.  The induced
splitting
\begin{equation}
    H^{p,q}(X)=H^{p,q}_{+}(X)\oplus H^{p,q}_{-}(X)
\end{equation}
selects the fields kept by the orientifold projection: in the
O3/O7 setting used here, the K\"ahler moduli are counted by
$h^{1,1}_{+}$, and the complex structure moduli in the flux
superpotential by $h^{2,1}_{-}$.  Sign conventions for
$\sigma^*\Omega$ differ in the literature; the convention displayed
above matches the one used for the flux formulae in these notes and
is the convention of \cite{Giddings:2001yu}, while
\cite{Grimm:2004uq} uses the opposite sign convention.
\end{definition}

For the toric hypersurfaces used in these lectures there is a
particularly effective way of finding orientifold involutions that
are inherited from the ambient toric variety.  Let
$X=\{f=0\}\subset V_\Sigma$ be a Batyrev hypersurface, with
$V_\Sigma$ described by Cox coordinates $x_p$ indexed by the rays
of $\Sigma$, hence by non-zero lattice points
$p\in\Delta^\circ\cap N$ in an MPCP subdivision.  The method of
\cite{Moritz:2023jdb} first replaces the problem of finding
holomorphic involutions of $X$ by the combinatorial problem of
finding involutions of the ambient Cox quotient.  After gauge fixing
in the ambient automorphism group, the representatives used in this
enumeration have the form
\begin{equation}
    \widetilde\sigma=\phi_{[t]}\circ L\, ,
\end{equation}
where $L$ is a lattice automorphism of $\Delta^\circ$ preserving the
relevant fan or regular subdivision, and $\phi_{[t]}$ is the torus
automorphism determined, after gauge fixing, by a half-integral
class $[t]\in \frac{1}{2}N/N$.  The condition
$\widetilde\sigma^2=1$ and the residual conjugacy action give a
finite list of representatives.  One then asks whether the locus of
$L$-symmetric height functions meets the secondary fan in a fine
regular star triangulation, or at least in a fine regular star
subdivision whose restrictions to all two-faces of
$\Delta^\circ$ are fine and regular.  This is the toric condition
under which the invariant ambient toric variety may be singular in
the symmetric limit while the generic invariant hypersurface remains
smooth.

The hypersurface equation is not arbitrary after the involution has
been chosen.  Writing
\begin{equation}
    f=\sum_{q\in\Delta\cap M}\psi_q\,s_q\, ,
\end{equation}
one imposes covariance
\begin{equation}
    \widetilde\sigma^*f=\gamma_f f
\end{equation}
for a fixed sign character $\gamma_f$.  Equivalently, the
coefficients $\psi_q$ are constrained along $L$-orbits, with phases
determined by $t$.  The sign of the induced action on the
Poincar\'e residue of $\Omega$ then distinguishes the O3/O7 case
from the O5/O9 case.  The fixed locus is computed directly from the
toric data.  If $\Sigma$ is the chosen invariant fan, its relevant
irreducible components are labelled by pairs $(\nu,\sigma)$ with
\begin{equation}
    \sigma\in\Sigma\ \text{pointwise fixed by }L,\qquad
    \nu\in H^-_L,\qquad
    t+\nu+\frac12\sum_{p\in\sigma(1)}p\in N\, ,
\end{equation}
where $H^-_L=P^-_L(N)/(2P^-_L(N))$ and
$P^-_L=(1-L)/2$ records the inequivalent anti-invariant toric
phases; in the displayed integrality condition $t$ and $\nu$ denote
chosen representatives of their classes.  One then removes
components contained in larger fixed components.  These ambient
fixed components are intersected with the invariant hypersurface.
Complex codimension-one components on $X$ give O7 divisors and
complex codimension-three components give O3 points.

This construction also computes the orientifold Hodge splitting.  In
the favourable case, the induced permutation of toric divisor
classes gives the $\pm$ eigenspaces of $H^{1,1}(X)$; the remaining
split of complex structure moduli is then obtained from the
Lefschetz fixed-point formula, using the Euler
characteristic of the fixed locus.  In the non-favourable case one
must enlarge this divisor calculation by the irreducible components
of toric divisors that split when restricted to $X$, notably those
coming from points interior to two-faces of $\Delta^\circ$.  This is
the ``orientifold weighted Hodge number'' computation of
\cite{Moritz:2023jdb}.  For the present lecture, the important
point is that the orientifold data become finite combinatorial input
read from the same toric description that also supplies the periods,
intersection numbers, and curve data.



For the rest of the lecture we mostly work with orientifold examples
satisfying
\begin{equation}
    h^{1,1}_{-}=0\, ,
    \qquad
    h^{2,1}_{+}=0\, .
\end{equation}
Then $h^{1,1}_{+}=h^{1,1}$ and $h^{2,1}_{-}=h^{2,1}$.  From this
point on we suppress the $\pm$ labels on the surviving cohomology
groups and Hodge numbers, unless the orientifold parity itself is the
point under discussion.  In particular, the flux lattice will be
written simply as $H^3(X,\mathbb Z)$ in the formulae below.






\subsection{Flux quantisation, periods, and the GVW input}
\label{ssec:L2-flux-period-setup}\label{ssec:L2-periods}

We now isolate the computational input that produces $W_{\rm flux}$.
A flux configuration is not a real vector of continuous parameters. It
is a pair of integral cohomology classes
\begin{equation}\label{eq:flux-quant}
    [F_3]\in H^3(X,\mathbb Z)\, ,
    \qquad
    [H_3]\in H^3(X,\mathbb Z)\, .
\end{equation}
The orientifold projection has already selected the relevant surviving
lattice; by the convention fixed above we do not write parity labels
in the rest of the lecture.  Integrality is the usual Dirac
quantisation condition: in string units, the periods over integral
three-cycles are integers.  Large gauge transformations of the
underlying potentials $C_2$ and $B_2$ identify potentials whose
periods differ by integers, so the cohomology class is the invariant
datum.

Choose a symplectic basis $\{\alpha_A,\beta^A\}$ of the relevant
integral cohomology lattice, with $A=0,\ldots,h^{2,1}$ in the
unprojected notation and with the corresponding surviving range after
the orientifold projection.  The pairing is
\begin{equation}\label{eq:L2-symplectic-pairing}
    \int_X\alpha_A\wedge\beta^B
    =
    \delta_A^B\, ,
    \qquad
    \int_X\alpha_A\wedge\alpha_B
    =
    \int_X\beta^A\wedge\beta^B
    =
    0\, .
\end{equation}
The periods of the holomorphic three-form are the coefficients in
\begin{equation}\label{eq:periods}
    \Omega(z)
    =
    X^A(z)\,\alpha_A-\mathcal F_A(z)\,\beta^A\, .
\end{equation}
Near a generic point of $\mathcal M_{\rm cs}(X)$, $h^{3,0}(X)=1$
and Griffiths transversality make the $X^A$ projective coordinates on
complex structure moduli space, so
$\mathcal F_A=\partial\mathcal F/\partial X^A$ for a homogeneous
degree-two prepotential $\mathcal F$ \cite{Candelas:1990rm}.  The
same period vector also gives the complex structure K\"ahler
potential,
\begin{equation}\label{eq:K-cs}
    K_{\rm cs}
    =
    -\log\!\Bigl(-\mathrm{i}\int_X\Omega\wedge\bar\Omega\Bigr)
    =
    -\log\!\Bigl(\mathrm{i}\bigl(X^A\bar{\mathcal F}_A
    -\bar X^A\mathcal F_A\bigr)\Bigr)\, .
\end{equation}
Thus computing periods is already enough to compute the metric part of
the flux equations of motion.

The physical flux combination is
\begin{equation}
    G_3
    =
    F_3-\tau H_3\, ,
\end{equation}
where $\tau$ is the axio-dilaton introduced in
\eqref{eq:axio-dilaton}.  The flux contribution to the
four-dimensional superpotential is the Gukov--Vafa--Witten period
pairing \cite{Gukov:1999ya},
\begin{equation}\label{eq:W-flux}
    W_{\rm flux}(\tau,z)
    =
    \int_X\Omega(z)\wedge G_3(\tau)\, .
\end{equation}
Writing
\begin{equation}
    F_3=f^A\alpha_A-f_A\beta^A\, ,
    \qquad
    H_3=h^A\alpha_A-h_A\beta^A\, ,
    \qquad
    f^A,f_A,h^A,h_A\in\mathbb Z\, ,
\end{equation}
one obtains, up to the sign convention for the symplectic pairing,
\begin{equation}\label{eq:L2-Wflux-period-vector}
    W_{\rm flux}
    =
    (f^A-\tau h^A)\,\mathcal F_A(z)
    -(f_A-\tau h_A)\,X^A(z)\, .
\end{equation}
This is the main reduction.  A choice of flux is an integer vector,
and the superpotential is the corresponding integer linear combination
of period functions.  Therefore the next mathematical task is not to
solve for vacua immediately, but to compute the period vector.

At large complex structure, the period vector splits into a polynomial
part and exponentially suppressed corrections.
\footnote{A \emph{large complex structure} (LCS) point is a boundary
point of $\mathcal M_{\rm cs}(X)$ at which the monodromy of the
periods of $\Omega$ is maximally unipotent.  Under mirror symmetry it
corresponds to the large volume limit on the mirror $\widetilde X$.}
Mirror symmetry computes these corrections from the large volume
prepotential on $\widetilde X$.  The coefficients of the instanton
part of that prepotential are genus-zero GV invariants.  This is the
precise reason the enumerative calculation below enters the lecture:
it supplies the instanton-corrected period entries used in
\eqref{eq:L2-Wflux-period-vector}.

The supersymmetric flux equations for the axio-dilaton and complex
structure moduli are
\begin{equation}\label{eq:L2-flux-Fterms}
    D_\tau W_{\rm flux}=0\, ,
    \qquad
    D_a W_{\rm flux}=0\, .
\end{equation}
The ten-dimensional GKP interpretation of these equations is important
background, but it is not needed for the period computation itself; see
\cite{Giddings:2001yu} and \cite{McAllister:2025qwq}, \S\,3.3.2.


% ------------------------------------------------------------------

\subsection{Computational mirror symmetry and Gopakumar--Vafa invariants}\label{ssec:L2-GV}



We have expressed $W_{\rm flux}$ as an integral linear combination of
periods.  The next task is therefore to compute those periods.  For
Calabi--Yau hypersurfaces in toric varieties this is most efficiently
done by mirror symmetry.  Periods on the Type IIB Calabi--Yau $X$ are
matched to the large volume $A$-model prepotential on the mirror
$\widetilde X$.  The exponentially small terms in that prepotential
are organised by genus-zero Gopakumar--Vafa invariants of curve
classes on $\widetilde X$ \cite{Gopakumar:1998ii,Gopakumar:1998jq}.
The computational mirror-symmetry programme of
\cite{Demirtas:2023als,RiosTascon:2023thesis} turns this statement
into an algorithm that remains practical for large Picard rank and for
non-simplicial Mori cones.

For this workshop audience it is useful to separate the conceptual
content from the algorithmic advance.  The conceptual framework is the
classical one: mirror symmetry near large complex structure
\cite{Candelas:1990rm,Hosono:1993qy}, Batyrev's toric construction
\cite{Batyrev:1993oya}, and the GV integral reorganisation of
Gromov--Witten theory \cite{Gopakumar:1998ii,Gopakumar:1998jq}.  The
newer contribution of \cite{Demirtas:2023als} is algorithmic: it gives
a practical way to compute the truncated instanton expansion when
$h^{1,1}$ is large and the relevant cones are high-dimensional.

Following the flow of \cite{Demirtas:2023als}, but using the notation
of the present lectures, the lecture-facing story is:
\begin{enumerate}
    \item compute the periods of $\Omega$ near an LCS point of
    $\mathcal M_{\rm cs}(X)$;
    \item identify the same period vector, by mirror symmetry, with
    the Type IIA large volume period vector on $\widetilde X$;
    \item read off the instanton part of the $A$-model prepotential;
    \item use the multiple-cover relation to extract the genus-zero GV
    integers;
    \item use toric GKZ/Frobenius series to carry this out without
    constructing a global Picard--Fuchs basis.
\end{enumerate}

\begin{remark}[Notation comparison with related references]\label{rem:L2-CMS-notation}
Throughout these lectures $X$ denotes the Calabi--Yau threefold used
for the Type IIB compactification, and $\widetilde X$ denotes its
mirror.  Thus Type IIA in the mirror symmetry comparison lives on
$\widetilde X$.  This is the convention used in the TASI discussion
of computational mirror symmetry \cite{McAllister:2025qwq},
\S5.5.2.2.  Reference~\cite{Demirtas:2023als} uses the opposite
labelling in parts of its exposition; the formulae below have been
translated into the notation of these lectures (hopefully consistently).
\end{remark}

\subsubsection*{The $A$-model prepotential and period vector}

Fix the mirror Calabi--Yau threefold $\widetilde X$.  Choose integral
curve classes $\{\widetilde\gamma^{a}\}\subset
H_2(\widetilde X,\mathbb Z)$, and choose a dual basis
$\{\widetilde\omega_a\}\subset H^{1,1}(\widetilde X)$ satisfying
\begin{equation}
    \int_{\widetilde\gamma^{a}}\widetilde\omega_b=\delta^{a}_{b}\, .
\end{equation}
We introduce complex coordinates
\begin{equation}\label{eq:L2-A-model-coord}
    \mathfrak t^a
    =
    \int_{\widetilde\gamma^{a}}(\widetilde B+
    \mathrm{i}\widetilde J)
    =
    b^a+\mathrm{i}t^a\, .
\end{equation}
Here $\widetilde J$ is the K\"ahler form on $\widetilde X$ and
$\widetilde B$ is the mirror-side $B$-field.  This makes
$\mathfrak t^a$ the Type IIA complexified Kähler parameters on $\widetilde X$.

In the large volume chamber, the genus-zero $A$-model prepotential has
the form
\begin{equation}
    F_0(\mathfrak t)
    =
    F_{0,{\rm poly}}(\mathfrak t)
    +
    F_{0,{\rm inst}}(\mathfrak t)\, .
\end{equation}
The polynomial part is fixed by classical topology: the triple
intersections and second-Chern-class integrals
\begin{equation}
    \widetilde\kappa_{abc}
    =
    \int_{\widetilde X}
    \widetilde\omega_a\wedge\widetilde\omega_b\wedge
    \widetilde\omega_c\, ,
    \qquad
    \widetilde c_a
    =
    \int_{\widetilde X}c_2(\widetilde X)\wedge
    \widetilde\omega_a\, ,
\end{equation}
together with the Euler characteristic $\chi(\widetilde X)$.
In the conventions of \cite{Demirtas:2023als}, translated to the
present notation, this polynomial part is
\begin{equation}\label{eq:L2-F-poly-CMS}
    F_{0,{\rm poly}}(\mathfrak t)
    =
    -\frac{1}{6}\widetilde\kappa_{abc}
      \mathfrak t^a\mathfrak t^b\mathfrak t^c
    +\frac{1}{2}\widetilde a_{ab}\mathfrak t^a\mathfrak t^b
    +\frac{1}{24}\widetilde c_a\mathfrak t^a
    +\frac{\zeta(3)\chi(\widetilde X)}{2(2\pi\mathrm{i})^3}\, .
\end{equation}
The quadratic coefficients $\widetilde a_{ab}$ record the chosen
integral symplectic basis.  In the convention used by
\cite{Demirtas:2023als} and \cite{McAllister:2025qwq}, \S5.5.2, they
are taken to be
\begin{equation}\label{eq:L2-a-ab}
    \widetilde a_{ab}
    =
    \frac{1}{2}
    \begin{cases}
        \widetilde\kappa_{aab}\, , & a\geq b\, ,\\
        \widetilde\kappa_{abb}\, , & a<b\, .
    \end{cases}
\end{equation}
Polynomial signs and lower-degree terms vary with period-vector
conventions; \eqref{eq:L2-F-poly-CMS} is the convention used for the
period matching below.

Let $\MC(\widetilde X)$ denote the Mori cone, i.e.\ the cone of
effective curve classes on $\widetilde X$.  In the chosen curve basis,
a class is written
$\widetilde{\mathbf q}=\widetilde q_a\widetilde\gamma^{a}$, and
$\widetilde{\mathbf q}\cdot\mathfrak t$ means
$\widetilde q_a\mathfrak t^a$.  The instanton part is most
transparently written first in terms of the genus-zero
Gromov--Witten invariants of $\widetilde X$:
\begin{equation}\label{eq:F-inst-GW}
    F_{0,{\rm inst}}(\mathfrak t)
    =
    -\frac{1}{(2\pi\mathrm{i})^3}
    \sum_{0\ne\widetilde{\mathbf q}
          \in\MC(\widetilde X)\cap H_2(\widetilde X,\mathbb Z)}
    N_{\widetilde{\mathbf q}}\,
    q^{\widetilde{\mathbf q}}\, ,
    \qquad
    q^{\widetilde{\mathbf q}}
    =
    \mathrm{e}^{2\pi\mathrm{i}\,
      \widetilde{\mathbf q}\cdot\mathfrak t}\, .
\end{equation}
Here $N_{\widetilde{\mathbf q}}\in\mathbb Q$ are genus-zero
Gromov--Witten invariants.  The GV integers are the coefficients that
repackage this same series into the polylogarithmic form
\begin{equation}\label{eq:F-inst}
    F_{0,{\rm inst}}(\mathfrak t)
    =
    -\frac{1}{(2\pi\mathrm{i})^3}
    \sum_{0\ne\widetilde{\mathbf q}
          \in\MC(\widetilde X)\cap H_2(\widetilde X,\mathbb Z)}
    n_{\widetilde{\mathbf q}}\,
    \mathrm{Li}_3\!\left(
      \mathrm{e}^{2\pi\mathrm{i}\,
      \widetilde{\mathbf q}\cdot\mathfrak t}
    \right)\, .
\end{equation}
This is the exact location at which the integers enter the period
calculation: they are the coefficients of the exponentially suppressed
part of the mirror prepotential, and hence of the corresponding
period vector.

Before continuing with the toric computation, we spell out what these
integers mean.  Given genus-zero Gromov--Witten invariants
$N_{\widetilde{\mathbf q}}$, the genus-zero GV coefficients are
defined by the multiple-cover relation
\begin{equation}\label{eq:L2-mult-cover}
    \sum_{0\ne\widetilde{\mathbf q}}
    N_{\widetilde{\mathbf q}}\,q^{\widetilde{\mathbf q}}
    =
    \sum_{0\ne\widetilde{\mathbf q}}
    n_{\widetilde{\mathbf q}}
    \sum_{k\geq1}\frac{1}{k^3}
    q^{k\widetilde{\mathbf q}}\, .
\end{equation}
Equivalently,
$N_{\widetilde{\mathbf q}}=\sum_{k\mid\widetilde{\mathbf q}}
k^{-3}n_{\widetilde{\mathbf q}/k}$.  Physically,
$n_{\widetilde{\mathbf q}}$ is a genus-zero BPS index obtained from
M2-branes wrapping holomorphic curves in class
$\widetilde{\mathbf q}$ in M-theory on $\widetilde X$; mathematically,
it is the integral reorganisation of rational GW invariants predicted
by Gopakumar--Vafa theory.  In computations, integrality of the
extracted $n_{\widetilde{\mathbf q}}$ is therefore a non-trivial
consistency check, not a tautology of the GW definition.

\begin{example}[Quintic threefold]\label{ex:L2-quintic-GV}
For the quintic threefold the first few genus-zero GV invariants are
\begin{equation}
    n_1=2875\, ,
    \qquad
    n_2=609\,250\, ,
    \qquad
    n_3=317\,206\,375\, ,
    \qquad \ldots\, .
\end{equation}
The number $n_1=2875$ agrees with the classical count of lines on a
generic quintic.  Higher-degree invariants must be read as virtual,
deformation-invariant counts after subtracting multiple-cover
contributions, not as naive counts of smooth embedded rational curves
\cite{Candelas:1990rm}.  For instance,
\begin{equation}
    N_2
    =
    n_2+\frac{n_1}{2^3}
    =
    609\,250+\frac{2875}{8}
    =
    609\,609.375\, .
\end{equation}
The rational GW coefficient has been repackaged as an integer GV count
plus the multiple-cover contribution of degree-one curves.
\end{example}

The corresponding projective period vector is obtained from the
homogeneous prepotential
$F(X^0,X^a)=(X^0)^2F_0(X^a/X^0)$.  In the gauge $X^0=1$ and
$X^a=\mathfrak t^a$, let
$F_a=\partial F_0/\partial\mathfrak t^a$.  Euler homogeneity gives
$F_0^{\rm hom}\coloneqq \partial F/\partial X^0=2F_0-\mathfrak t^aF_a$ in
this gauge, so the $A$-model period vector can be written as
\begin{equation}\label{eq:L2-Amodel-period-vector}
    \Pi_A(\mathfrak t)
    =
    \begin{pmatrix}
        2F_0(\mathfrak t)-\mathfrak t^aF_a(\mathfrak t)\\
        F_a(\mathfrak t)\\
        1\\
        \mathfrak t^a
    \end{pmatrix}
\end{equation}
up to the chosen symplectic-basis convention.  The Frobenius
calculation below produces the $B$-model representative of this same
vector after applying the mirror map.  Thus the comparison to keep in
mind is
\begin{equation}
    \Pi_B(\psi)
    \quad\stackrel{\rm mirror}{=}\quad
    \Pi_A(\mathfrak t(\psi))\, .
\end{equation}

\subsubsection*{Mirror computation for toric hypersurfaces}

The direct $A$-model definition of the integers on $\widetilde X$
is not the computational route used here.  Instead one computes the
Type IIB $B$-model periods on $X$.  This is the classical side of
the mirror pair: the relevant vector-multiplet geometry is encoded
by periods of the holomorphic three-form $\Omega$ on $X$.  Mirror
symmetry says that, after choosing the correct integral symplectic
basis and normalisation, these periods reproduce the Type IIA
$A$-model prepotential of $\widetilde X$.

For Batyrev hypersurfaces, the mirror pair is constructed from a
dual pair of four-dimensional reflexive polytopes
$(\Delta^\circ,\Delta)$.  A fine, regular, star triangulation of
$\Delta^\circ$ defines a toric fourfold $V$; in the usual logic \cite{Batyrev:1993oya} the generic anti-canonical hypersurface
$X\subset V$ is a smooth Calabi--Yau threefold for the Type IIB
compactification.  Exchanging $\Delta^\circ$ and $\Delta$ gives the
mirror $\widetilde X\subset\widetilde V$, whose Type IIA
large volume prepotential contains the GV invariants we want.  The
period integral below is taken on $X$, but the coefficient formula is
naturally indexed by ambient curve-class data of the mirror toric
variety $\widetilde V$.  

The lattice relations among the monomials in the defining polynomial
of $X$ become the charge data entering the LCS coordinates.  We use
capital indices $I,J$ for toric monomials and ambient toric divisors,
and lowercase indices $a,b,c$ for complex structure coordinates on
$X$, equivalently for mirror curve/divisor bases on $\widetilde X$.
Write
\begin{equation}\label{eq:L2-defining-polynomial}
    f \;=\; \Psi^0\,S_0 - \sum_I \Psi^I\,S_I
\end{equation}
for the defining polynomial of $X\subset V$. The
$S_0,\,S_I$ are the monomial sections of the anticanonical bundle
$-K_{V}$ on the original-side toric fourfold $V$; they are indexed by the lattice
points of the dual polytope $\Delta$ (the same polytope whose fan
defines the mirror $\widetilde V$), with $S_0$ the interior-point
monomial corresponding to the origin. The
$\Psi^0,\Psi^I\in\mathbb{C}$ are their complex coefficients; only
the ratios $\Psi^I/\Psi^0$ enter the complex structure moduli data.
The integer charge matrix $Q^a_I$ ($a=1,\ldots,h^{1,1}(\widetilde
V)$ in the ambient toric basis, restricted afterwards to the
hypersurface directions used above) encodes the linear relations
$\sum_I Q^a_I\,\nu_I=0$ among the lattice-point vectors
$\nu_I\in\Delta$; equivalently, it is the matrix of Mori-cone
generators on the ambient mirror toric fourfold $\widetilde V$, and
its rows are the relations among the monomials $S_I$ on $V$.
Multiplicative relations among the monomials then give local
gauge-invariant complex structure coordinates
\begin{equation}
    \psi^a=\prod_I
      \left(\frac{\Psi^I}{\Psi^0}\right)^{Q^a_I}\, .
\end{equation}
Interior-facet points are treated separately: the corresponding
ambient divisors do not meet a generic hypersurface, and part of the
toric automorphism group can be used to remove redundant polynomial
coefficients.  This is one reason why the toric input must be handled
as a quotient construction rather than as a list of monomials with
independent coefficients.

Thus the Batyrev input plays three simultaneous roles in this
mirror-period computation.  It gives local complex structure
coordinates on $X$; it gives divisor, intersection, and Chern-class
data on the mirror $\widetilde X$; and it gives the charge matrix
$Q^a_I$ that determines the GKZ/Frobenius period computation.  This is the reason
the method is useful for these lectures: the same finite lattice data
feed both the Type IIB period calculation and the Type IIA
curve-counting calculation.

The basic period is obtained by restricting to the dense algebraic
torus $(\mathbb{C}^*)^4\subset V$ in two steps. First, the
Poincar\'e residue defines the holomorphic three-form on $X$:
\begin{equation}
    \Omega
      =
      \oint_{f=0}\!\left[\mathrm{d}y^1\wedge\cdots\wedge\mathrm{d}y^4
      \,
      \frac{1}{(2\pi\mathrm{i})^4 f(y)}\right]\, .
\end{equation}
Second, the fundamental period $\varpi_0(\psi)$ is the
integral of $\Omega$ over the standard small-torus cycle on
$(\mathbb{C}^*)^4\subset V$ at $\psi=0$; equivalently, it is the
constant-term Fourier mode of the right-hand side above.
The resulting fundamental period, expressed as a sum over integer
effective ambient curve classes in the chosen toric charge basis, is
\begin{equation}\label{eq:L2-fundamental-period}
    \varpi_0(\psi) = \Psi^{0}\int_{T^{3}}\, \Omega
    =
    \sum_{\widetilde{\mathbf q}\in\MC(\widetilde V)
          \cap H_2(\widetilde V,\mathbb Z)}
    c_{\widetilde{\mathbf q}}\,
    \psi^{\widetilde{\mathbf q}}\, ,
\end{equation}
where $T^{3}\subset X$ defines a certain solution branch for $f=0$ (the SYZ fiber \cite{Strominger:1996it})
and where
\begin{equation}
c_{\widetilde{\mathbf q}}
    \coloneqq 
    \frac{\Gamma(1+\widetilde q_aQ^a_0)}
         {\prod_I\Gamma(1+\widetilde q_aQ^a_I)}\, ,
    \qquad
    \psi^{\widetilde{\mathbf q}}
    \coloneqq 
    \prod_a(\psi^a)^{\widetilde q_a}\, .
\end{equation}

Here $Q^a_I$ ($I\ne 0$) is the integer charge matrix
introduced in \eqref{eq:L2-defining-polynomial}, and
$Q^a_0$ is the charge attached to the interior monomial $S_0$ fixed
by the Calabi--Yau condition on the toric data.
The coefficient $c_{\widetilde{\mathbf q}}$ is the basic
$\Gamma$-function ratio that all the other period-series coefficients
below are built from.  Its analytic continuation in
$\widetilde{\mathbf q}$, obtained by replacing the integer vector
$\widetilde{\mathbf q}$ with $\widetilde{\mathbf q}+\vec\rho$ for a
formal auxiliary vector $\vec\rho=(\rho_a)$, is what allows the
Frobenius derivatives below to be taken.

Initially this expression arises as a constant-term computation over
movable ambient curve classes.  It can be extended to the ambient
Mori cone $\MC(\widetilde V)$: outside the movable cone, denominator
poles in the gamma expression kill the forbidden terms
\cite{Demirtas:2023als}.  The important point is that the period is
computed as a concrete toric hypergeometric series.
The notation deliberately distinguishes the ambient cone
$\MC(\widetilde V)$ behind \eqref{eq:L2-fundamental-period} from the
hypersurface cone $\MC(\widetilde X)$ in \eqref{eq:F-inst}. The
former is a computational lattice for the hypergeometric period; the
latter is the geometric cone in which the enumerative invariants
live. The mirror-period method of \cite{Demirtas:2023als} works
because these two pieces of information can be compared in a common
toric charge basis, with the caveats about $\MC_{\widetilde V}^{\cap}$
stated below.

The full Picard--Fuchs system is the differential system satisfied
by all periods.  In the toric setting it is obtained from the
GKZ system determined by the same lattice relations $Q^a_I$.
Factoring the GKZ system down to the true Picard--Fuchs system is
usually expensive.  The HKTY method avoids this as follows
\cite{Hosono:1993qy}.  At a large complex structure point, the
monodromy is maximally unipotent and is mirror to the large volume
$B$-field monodromy of Type IIA on $\widetilde X$.  Hence the
logarithmic Frobenius solutions, obtained by differentiating the
term $c_{\widetilde{\mathbf q}+\vec\rho}\,
\psi^{\widetilde{\mathbf q}+\vec\rho}$ of
\eqref{eq:L2-fundamental-period} with respect to the auxiliary
variables $\rho_a$ before setting $\vec\rho=0$, give the needed local
period basis:
\begin{align}
    \varpi^a(\psi)
      &=
      \sum_{\widetilde{\mathbf q}}
      \left.
      \frac{\partial_{\rho_a}}{2\pi\mathrm{i}}
      \Bigl(c_{\widetilde{\mathbf q}+\vec\rho}\,
            \psi^{\widetilde{\mathbf q}+\vec\rho}\Bigr)
      \right|_{\vec\rho=0},\\
    \varpi^{ab}(\psi)
      &=
      \sum_{\widetilde{\mathbf q}}
      \left.
      \frac{\partial_{\rho_a}\partial_{\rho_b}}
           {(2\pi\mathrm{i})^2}
      \Bigl(c_{\widetilde{\mathbf q}+\vec\rho}\,
            \psi^{\widetilde{\mathbf q}+\vec\rho}\Bigr)
      \right|_{\vec\rho=0},\\
    \varpi^{abc}(\psi)
      &=
      \sum_{\widetilde{\mathbf q}}
      \left.
      \frac{\partial_{\rho_a}\partial_{\rho_b}\partial_{\rho_c}}
           {(2\pi\mathrm{i})^3}
      \Bigl(c_{\widetilde{\mathbf q}+\vec\rho}\,
            \psi^{\widetilde{\mathbf q}+\vec\rho}\Bigr)
      \right|_{\vec\rho=0}\, .
\end{align}

\medskip
\noindent It is useful to record once and for all the coefficient
data implicit in these Frobenius solutions.  Define
\begin{equation}\label{eq:L2-c-coefs}
    c^a_{\widetilde{\mathbf q}}
      \coloneqq 
      \left.\partial_{\rho_a}
            c_{\widetilde{\mathbf q}+\vec\rho}\right|_{\vec\rho=0},
    \qquad
    c^{ab}_{\widetilde{\mathbf q}}
      \coloneqq 
      \left.\partial_{\rho_a}\partial_{\rho_b}
            c_{\widetilde{\mathbf q}+\vec\rho}\right|_{\vec\rho=0},
\end{equation}
and assemble the formal power series
\begin{equation}\label{eq:L2-c-series}
    c^a(\psi)
      \coloneqq 
      \sum_{\widetilde{\mathbf q}\in\MC(\widetilde V)}
      c^a_{\widetilde{\mathbf q}}\,\psi^{\widetilde{\mathbf q}},
    \qquad
    c^{ab}(\psi)
      \coloneqq 
      \sum_{\widetilde{\mathbf q}\in\MC(\widetilde V)}
      c^{ab}_{\widetilde{\mathbf q}}\,\psi^{\widetilde{\mathbf q}}.
\end{equation}
These are the gamma-function-derived power series referenced
repeatedly in the period-matching formulae below.  Algorithmically,
$c_{\widetilde{\mathbf q}}$, $c^a_{\widetilde{\mathbf q}}$ and
$c^{ab}_{\widetilde{\mathbf q}}$ are independent and can be computed
in parallel; cf.\ \cite{Demirtas:2023als}\,\S5.


Matching the large complex structure monodromies to the large volume
monodromies fixes the constant change of basis to an integral
symplectic period vector.  The comparison is with the $A$-model vector
\eqref{eq:L2-Amodel-period-vector}: the first two blocks correspond
to $2F_0-\mathfrak t^aF_a$ and $F_a$, while the last two blocks
correspond to $1$ and the flat coordinates $\mathfrak t^a$.  The
classical data $\widetilde\kappa_{abc}$, $\widetilde c_a$, and
$\chi(\widetilde X)$ fix the polynomial part;
the remaining power series is the instanton contribution.

In the period-basis convention adopted in \cite{Demirtas:2023als},
the $B$-model representative of \eqref{eq:L2-Amodel-period-vector} is
\begin{equation}\label{eq:L2-CMS-period-vector}
    \Pi(\psi)
      =
      \frac{1}{\varpi_0}
      \begin{pmatrix}
      \frac{1}{3!}\widetilde\kappa_{abc}\varpi^{abc}
        +\frac{1}{12}\widetilde c_a\varpi^a\\
      -\frac{1}{2}\widetilde\kappa_{abc}\varpi^{bc}
        +\widetilde a_{ab}\varpi^b\\
      \varpi_0\\
      \varpi^a
      \end{pmatrix}\, .
\end{equation}
The ordering of components is a convention, but the lesson is
coordinate-free: the polynomial part of the period vector is fixed
by classical topology, and the remaining power series is where the
instanton/GV data appear.

\subsubsection*{Extracting GV invariants from the periods}

The flat coordinate is
\begin{equation}\label{eq:L2-mirror-map}
    \mathfrak{t}^a
      =
      \frac{\varpi^a}{\varpi_0}
      =
      \frac{\log\psi^a}{2\pi\mathrm{i}}
      +\frac{1}{2\pi\mathrm{i}}\,
       \frac{c^a(\psi)}{\varpi_0(\psi)}\, ,
\end{equation}
which is the \emph{mirror map}: it identifies the $B$-model
algebraic complex structure coordinates $\psi^a$ on $X$ with the
$A$-model complexified K\"ahler coordinates $\mathfrak t^a$ on
$\widetilde X$ originally defined via $\mathfrak t^a =
\int_{\widetilde\gamma^{a}}(\widetilde B+\mathrm{i}\widetilde J)$ in
\eqref{eq:L2-A-model-coord} above.  Thus
\begin{equation}
    q^{\widetilde{\mathbf q}}
      =
      \mathrm{e}^{2\pi\mathrm{i}\,
        \widetilde{\mathbf q}\cdot\mathfrak{t}}
      =
      \psi^{\widetilde{\mathbf q}}
      \exp\!\left(
        \widetilde q_a\,\frac{c^a(\psi)}{\varpi_0(\psi)}
      \right)\, .
\end{equation}
Differentiating $F_{0,{\rm inst}}$ in
\eqref{eq:F-inst} and comparing with the period component
$F_a=\partial F_0/\partial\mathfrak{t}^a$ gives the master
identity
\begin{equation}\label{eq:L2-CMS-GV-master}
    \sum_{\widetilde{\mathbf q}\in\MC(\widetilde X)}
    \widetilde q_a\,n_{\widetilde{\mathbf q}}
    \mathrm{Li}_2\!\left(
      \psi^{\widetilde{\mathbf q}}
      \exp\!\left(
        \widetilde q_b\,\frac{c^b(\psi)}{\varpi_0(\psi)}
      \right)
    \right)
    =
    \frac{1}{2}\widetilde\kappa_{abc}
    \frac{\widehat c^{bc}(\psi)-c^b(\psi)c^c(\psi)}
         {\varpi_0(\psi)}\, .
\end{equation}
Here $c^a(\psi)$ is the series defined in
\eqref{eq:L2-c-series}.
The hatted series $\widehat c^{ab}(\psi)$ is the same Frobenius
coefficient $c^{ab}(\psi)$ from \eqref{eq:L2-c-series}, shifted by an
explicit multiple of the fundamental period:
\begin{equation}\label{eq:L2-c-hat}
    \widehat c^{ab}(\psi)
      \coloneqq 
      c^{ab}(\psi)
      -\frac{\pi^2}{6}\Bigl(Q^a_0Q^b_0
        -\sum_I Q^a_IQ^b_I\Bigr)\varpi_0(\psi)\,  .
\end{equation}
Equation~\eqref{eq:L2-CMS-GV-master} is the practical bridge from
periods to integers: expand both sides in powers of $\psi^a$, solve
degree by degree, and then convert the rational Gromov--Witten
coefficients into integral GV invariants using the multiple-cover
formula.

There is one important cone issue.  The period series is naturally
written using the ambient Mori cone, but the GV invariants of
interest live in the Mori cone of the hypersurface.  Moreover,
different triangulations of the mirror polytope $\Delta$ give
different ambient Mori cones of $\widetilde V$, while many ambient
birational transformations do not change the smooth mirror
hypersurface.  The computation therefore uses
$\MC_{\widetilde V}^\cap$: an intersection of ambient Mori cones over
phases for which the generic mirror hypersurface remains smooth.  It
is an outer approximation to $\MC(\widetilde X)$, not a theorem that
every point in it is an effective curve class on $\widetilde X$.
This distinction matters when interpreting zeros and truncations.

The finite truncation must also be closed under decompositions of
effective classes.  If a coefficient in
\eqref{eq:L2-CMS-GV-master} at class $\widetilde{\mathbf q}$ can
receive contributions from
$\widetilde{\mathbf q}=\widetilde{\mathbf q}'
+\widetilde{\mathbf q}''$, then the coefficients at
$\widetilde{\mathbf q}'$ and $\widetilde{\mathbf q}''$ must already
be in the truncated computation.  In the notation of
\cite{Demirtas:2023als},
one may use a causal diamond
\begin{equation}
    \lozenge_{\widetilde V}(\widetilde{\mathbf q})
      =
      \MC_{\widetilde V}^\cap
      \cap(\widetilde{\mathbf q}-\MC_{\widetilde V}^\cap)
\end{equation}
and require $\lozenge_{\widetilde V}(\widetilde{\mathbf q})$ to be
included for every class $\widetilde{\mathbf q}$ whose invariant is
extracted.  A common practical alternative is the degree method:
choose an integral grading vector $\mathbf d$ in the strict interior
of the dual K\"ahler cone and retain all classes with
$\mathbf d\cdot\widetilde{\mathbf q}\leq \ell$.  Since degree is
additive on the semigroup, this truncation is automatically closed
under decompositions.  Its cost scales like
$\ell^{h^{1,1}}$, so at large $h^{1,1}$ one often computes
modest-degree bulk data together with much deeper data on selected
rays or low-dimensional faces.

The
implementation of \cite{Demirtas:2023als}  stores truncated power series sparsely, orders
monomials by degree, and uses hash maps to make polynomial
multiplication feasible in many variables.  The coefficients
$c_{\widetilde{\mathbf q}}$,
$c^a_{\widetilde{\mathbf q}}$, and
$c^{ab}_{\widetilde{\mathbf q}}$ are independent and can
be computed in parallel.  The invariants are then extracted degree by
degree:
after the invariants at a fixed degree are read off, their
$\mathrm{Li}_2(q^{\widetilde{\mathbf q}})$ contributions are
subtracted before moving to the next degree.  This avoids constructing
or inverting the full Picard--Fuchs system and is the reason the
method scales beyond the older one- or few-parameter computations.

\subsubsection*{Aside: Toric curves and analytic checks}



We introduce the phrase ``toric curve'' only at this point, and it
needs a qualification.  A generic anticanonical hypersurface
$\widetilde X\subset\widetilde V_{\widetilde\Sigma}$ is not itself a
toric variety.  Let $\widehat{\widetilde D}_I$ denote a toric divisor
of the ambient toric fourfold $\widetilde V_{\widetilde\Sigma}$, and
let $\widetilde D_I=\widehat{\widetilde D}_I\cap\widetilde X$ denote
its restriction to the hypersurface when this restriction is non-empty.
Following the analytic checks of
\cite{Braun:2017nhi,Demirtas:2023als}, a toric curve class here means
a curve class on $\widetilde X$ that is visible from the ambient toric
presentation, typically a complete intersection
\begin{equation}
    \widetilde C_{IJ}=\widetilde D_I\cap \widetilde D_J
      =
      \widetilde X\cap \widehat{\widetilde D}_I
      \cap \widehat{\widetilde D}_J\, .
\end{equation}
Such curves are especially useful because some of their GV invariants
can be computed directly, without mirror symmetry, and can therefore
serve as checks of the period computation.

Let $\widetilde p_I,\widetilde p_J\in\Delta$ be the lattice points
corresponding to
$\widehat{\widetilde D}_I,\widehat{\widetilde D}_J$.  A
``one-face divisor'' or ``two-face divisor'' means that the
corresponding lattice point lies in the relative interior of a one- or
two-dimensional face of $\Delta$, respectively.  The complete
intersection $\widetilde C_{IJ}$ is non-empty in the simple
toric-curve analysis used here only when the segment joining
$\widetilde p_I$ and $\widetilde p_J$ is an edge of the chosen
triangulation and this edge lies in a common two-face $\Theta_2$ of
$\Delta$.  Let $g$ be the number of
interior lattice points
of the polar-dual one-face $\Theta_2^\circ\subset\Delta^\circ$
(by polar duality on the reflexive pair
$(\Delta,\Delta^\circ)$: a 2-face $\Theta_2\subset\Delta$ is dual to
a 1-face $\Theta_2^\circ\subset\Delta^\circ$ with
$\dim\Theta_2 + \dim\Theta_2^\circ = 3$),
equivalently the genus attached to $\Theta_2$.  On the open patch
associated with $\Theta_2$, the mirror hypersurface is a branched
cover of
$g+1$ copies of a non-compact toric Calabi--Yau threefold
$V_{\Theta_2}$.  Consequently
\begin{equation}
    \widetilde C_{IJ}
      =
      \coprod_{i=1}^{g+1} \widetilde C_{IJ}^{(i)},
    \qquad
    \widetilde C_{IJ}^{(i)}\simeq\mathbb{P}^1\, .
\end{equation}
If either $\widetilde D_I$ or $\widetilde D_J$ is a two-face divisor,
the components $\widetilde C_{IJ}^{(i)}$ lie in distinct numerical
classes.  Otherwise they represent the same numerical class with
multiplicity $g+1$
\cite{Braun:2017nhi,Demirtas:2023als}.

Under the local hypotheses used in
\cite{Braun:2017nhi,Demirtas:2023als} for one component, the normal bundle in
$\widetilde X$ has the form
\begin{equation}
    \mathcal N_{\widetilde C_{IJ}^{(i)}/\widetilde X}
      \simeq
      \mathcal{O}_{\mathbb{P}^1}(m)\oplus
      \mathcal{O}_{\mathbb{P}^1}(n),
    \qquad
    m=\frac{\widetilde\kappa_{IIJ}}{g+1},\quad
    n=\frac{\widetilde\kappa_{IJJ}}{g+1},\quad
    m+n=-2\, .
\end{equation}
Assume $m\geq n$.  For the smooth moduli spaces considered in
\cite{Demirtas:2023als}, the genus-zero GV invariant is computed by
\begin{equation}
    \mathrm{GV}^0=(-1)^{\dim\mathcal{M}}\chi(\mathcal{M})\, .
\end{equation}
If the singular locus produced by shrinking the curve is compact,
then the relevant moduli space is $\mathbb{P}^{m+1}$, giving
\begin{equation}
    (-1)^{m+1}(m+2)\, .
\end{equation}
Putting the component structure together gives, for the cases covered
by the compact-singular-locus assumption,
\begin{equation}
    \mathrm{GV}^{0}_{[\widetilde C_{IJ}^{(i)}]}
      =
      \begin{cases}
      (-1)^{m+1}(m+2),
        & \text{if either }\widetilde D_I\text{ or }\widetilde D_J\text{ is a two-face divisor},\\
      g+1,
        & \text{otherwise}\, .
      \end{cases}
\end{equation}
There is one remaining case, $m=0$ with $\widetilde D_J$
a one-face divisor. Then $\widetilde D_J$ is a
$\mathbb{P}^1$-fibration over a curve of genus $g'$, where
$g'$ equals the number of interior lattice points of the dual
two-face of $\Delta^\circ$ containing the relevant edge, and the
moduli space of the fibre class is that base curve. The
corresponding invariant is $2g'-2$.

These formulas are not a substitute for the full mirror computation:
they apply to a special class of ambient-visible rational curves.
Their value for the lectures is conceptual and practical.  They show
what kind of curve geometry the integers detect, and they provide
local checks on the GV dictionaries used in lecture 3 for flop navigation and wall crossing in the (torically) extended Kähler cone.
These toric curve classes include the
Mori-extremal classes whose volumes can shrink at chamber walls of
the secondary fan, so their GV invariants are the enumerative data
that survive into the flop transitions of Lecture 3 and into the
explicit constructions it feeds.


\subsection{Application: perturbatively flat vacua from periods and GV invariants}\label{ssec:L2-PFV-application}


This final section is the promised application of the period and
GV machinery developed above.  Up to this point, the lecture has explained
how toric mirror symmetry computes the period vector and how the GV
integers enter its instanton expansion.  We now insert that computed
period vector into the flux equations of motion.  The aim is not yet a
full vacuum-energy construction; it is to show, in one explicit class of
examples, how enumerative data become numerical input for a flux
superpotential.

\subsubsection{From periods to equations of motion: perturbatively flat vacua}\label{ssec:L2-smallW0}

We have now reached the first payoff of the period
calculation.  Near an LCS point, mirror symmetry gives an explicit
series expansion for the period vector $\Pi(z)$, including the
worldsheet-instanton corrections controlled by GV invariants.  Once
the flux integers are chosen, the flux superpotential is therefore a
computable holomorphic function.  Schematically, if $f$ and $h$
collect the integral $F_3$- and $H_3$-flux quanta, then
\begin{equation}
    W_{\rm flux}(\tau,z)
      =
      (f-\tau h)\cdot \Pi(z)\, ,
\end{equation}
where the dot denotes the symplectic pairing in the chosen integral
period basis.  The next problem is no longer to compute periods.  It
is to solve the equations of motion for the axio-dilaton and complex
structure moduli.

For constant scalar fields, the general equations of
motion are the critical-point equations for the scalar potential,
\begin{equation}
    \partial_{\Phi^I}V=0\, ,
\end{equation}
with $\Phi^I=(\tau,z^a,T_i)$.  In the supersymmetric flux sector
for $(\tau,z^a)$, it is enough to solve the $F$-term equations
\begin{equation}\label{eq:L2-flux-EOM}
    D_\tau W_{\rm flux}=0,\qquad
    D_a W_{\rm flux}=0\, ,
\end{equation}
with $D_IW=\partial_IW+(\partial_IK)W$.  These conditions are equivalent to $\star G_3=\mathrm{i}G_3$ with $\star$ corresponding to the Hodge-star on $X$ and $G_{3}=F_3-\tau H_3$.
For $I=a$, the relevant
$K$ is precisely the Weil--Petersson K\"ahler potential
\eqref{eq:K-cs-WP}, and the metric entering the scalar potential is
\eqref{eq:WP-cs-Hodge}.  Thus the period computation supplies both
the function $W_{\rm flux}$ and the geometry with respect to which
its covariant derivatives are taken.  The equations
\eqref{eq:L2-flux-EOM} are the four-dimensional equations of motion
restricted to the supersymmetric critical-point ansatz for
$(\tau,z^a)$.  For a generic flux vector,
\eqref{eq:L2-flux-EOM} is a coupled system of transcendental
equations, because the periods contain logarithms and instanton
series.  
They are Diophantine because the unknown flux quanta are required to be integers.
With an integral symplectic basis fixed, the fluxes are integer vectors, so the equations \eqref{eq:L2-flux-EOM} impose algebraic constraints on integer coefficients rather than on arbitrary real parameters. In the LCS regime, the leading polynomial part makes this especially explicit: one is asking for integer flux vectors satisfying a system of linear or polynomial relations, often together with a tadpole bound. This is difficult because integer solutions are sparse and discontinuous; unlike ordinary numerical minimisation, one cannot move infinitesimally in flux space, and the search grows rapidly with $b_3(X)$ and the allowed tadpole range.
In practice one therefore usually solves the system \eqref{eq:L2-flux-EOM} numerically.



Instead, we make use of the \emph{perturbatively flat
vacuum mechanism} of Demirtas, Kim, McAllister and Moritz
\cite{Demirtas:2019sip}. 
Such perturbatively flat vacua are special because part of the
problem becomes \emph{algebraic}.  The idea is to choose the fluxes so that
the polynomial LCS approximation to $W_{\rm flux}$ has a flat
solution locus on which both $W_{\rm flux}$ and its first derivatives
vanish.  The exponentially small instanton terms, whose coefficients
are the GV invariants computed above, then lift this flat direction.
This produces a computable one-variable problem and can leave an
exponentially small value for
\begin{equation}
    W_0=\langle W_{\rm flux}\rangle\, .
\end{equation}
Thus the construction should be read as a controlled
shortcut through the flux equations of motion: most flux choices lead
to a numerical root-finding problem, while the PFV flux choices first
solve the polynomial LCS equations exactly and then use the
instanton-corrected periods to stabilise the remaining direction.
This is why PFVs are useful for later:
small $|W_0|$ pushes the (SUSY) K\"ahler modulus minimum to large divisor
volume.  But the logical order in this lecture is simply
\begin{equation}
    \text{periods at LCS}
    \quad\Longrightarrow\quad
    \text{flux equations of motion}
    \quad\Longrightarrow\quad
    \text{PFVs as analytic foothold.}
\end{equation}


\subsubsection*{Aside: The tadpole constraint on flux choices*}

The Diophantine problem is also bounded.  In a compact
orientifold, three-form flux carries D3-brane charge and must fit
inside the global D3-charge budget.  With the localised O-plane and
D7-brane curvature contributions moved to the right-hand side, the
standard convention used in the GKP/F-theory formulation is:
\begin{proposition}[D3 tadpole condition]\label{prop:tadpole}
For a Type IIB O3/O7 compactification with F-theory lift
$Y_4$ (an elliptically fibred Calabi-Yau fourfold), the fluxes and net mobile D3-charge obey
\begin{equation}\label{eq:tadpole}
    N_{\mathrm{D3}}-N_{\overline{\mathrm{D3}}}
    +\frac12\int_X H_3\wedge F_3
    =
    \frac{\chi(Y_4)}{24}\, .
\end{equation}
Equivalently, the flux contribution
$N_{\rm flux}=\int_X H_3\wedge F_3$ is restricted by the finite
D3-charge budget after a choice of localised sources.  Signs vary
with the ordering convention for $F_3$ and $H_3$; the convention above
matches the use of $G_3=F_3-\tau H_3$ in these notes.  See
\cite{Sethi:1996es,Giddings:2001yu}.
\end{proposition}

\subsubsection*{The PFV computation in LCS coordinates}

We now spell out the calculation, because this is the first point in
the lectures where the period technology from \S\ref{ssec:L2-GV}
directly produces a vacuum ingredient.
Near an LCS point, choose flat
complex structure coordinates $z^a$, $a=1,\ldots,h^{2,1}(X)$.
These are the $B$-model counterparts of the mirror
complexified K\"ahler coordinates $\mathfrak t^a$ used in the
mirror-period discussion of \S\ref{ssec:L2-GV}.  In the symplectic
period basis adapted to the LCS prepotential (the basis
underlying the period vector $\Pi$ of
\eqref{eq:L2-CMS-period-vector}, with slot ordering
$(\mathcal F_0, \mathcal F_a, X^0, X^a)^{\top}$),
write the flux vectors as
\begin{equation}\label{eq:L2-PFV-flux-ansatz}
    \vec f=(R_0,R_a,0,M^a)^{\top},
    \qquad
    \vec h=(0,K_a,0,\mathbf 0)^{\top},
    \qquad
    R_0,R_a,M^a,K_a\in\mathbb Z\, .
\end{equation}
The bold $\mathbf 0$ in the magnetic slot of $\vec h$ is the
vector zero (matching the vector $M^a$ in the same slot of $\vec
f$), not a scalar zero.
Let $\widetilde\kappa_{abc}$, $\widetilde a_{ab}$, and
$\widetilde c_a$ denote the cubic, quadratic, and linear
polynomial data of the LCS prepotential in this basis.  Up to the
common overall normalisation convention for $W_{\rm flux}$, the
period expansion gives
\begin{equation}\label{eq:L2-PFV-W-split}
    W_{\rm flux}=W_{\rm poly}+W_{\rm exp}\, ,
\end{equation}
with
\begin{align}
    W_{\rm poly}(z,\tau)
    &=
      \frac12 M^a\widetilde\kappa_{abc}z^bz^c
      -\tau K_a z^a
      +(R_a-\widetilde a_{ab}M^b)z^a
      +\left(R_0-\frac{1}{24}M^a\widetilde c_a\right)\, ,
      \label{eq:L2-PFV-Wpoly}\\
    W_{\rm exp}(z)
    &=
      -\frac{1}{(2\pi)^2}
      \sum_{\widetilde{\mathbf q}\in\MC(\widetilde X)}
      n_{\widetilde{\mathbf q}}\,
      \widetilde q_a M^a\,
      \operatorname{Li}_2
      \!\left(
        \mathrm{e}^{2\pi\mathrm{i}\,\widetilde q_a z^a}
      \right)\, .
      \label{eq:L2-PFV-Wexp}
\end{align}
Here $\widetilde X$ is the mirror used to compute the periods, and
$n_{\widetilde{\mathbf q}}$ are the relevant genus-zero GV
invariants of curve classes on $\widetilde X$.  Thus the same
integers that entered the instanton part of the prepotential in
\eqref{eq:F-inst} now appear as coefficients in the flux
superpotential.

The point of the special flux ansatz is that the polynomial part can
be made to vanish before using the exponentially small terms.  If
\begin{equation}\label{eq:L2-PFV-integrality}
    R_a=\widetilde a_{ab}M^b\in\mathbb Z\, ,
    \qquad
    R_0=\frac{1}{24}M^a\widetilde c_a\in\mathbb Z\, ,
\end{equation}
then the constant and linear terms in $W_{\rm poly}$ are absent.
Since $\widetilde a_{ab}$ is in general half-integral
(see \eqref{eq:L2-a-ab}), the integrality conditions are best viewed
as congruence conditions on $M$.  Projecting to the $M$-variables,
the allowed choices are the kernel of the homomorphism
\begin{equation}
    \mathbb Z^{h^{2,1}(X)}
    \longrightarrow
    (\mathbb Q/\mathbb Z)^{h^{2,1}(X)+1},\qquad
    M\longmapsto
    \left(\widetilde a_{ab}M^b,\,
    \frac{1}{24}\widetilde c_aM^a\right)\bmod\mathbb Z\, .
\end{equation}
This kernel always contains $0$ and is a finite-index
sublattice of the integer flux lattice (the codomain is a finite
abelian group, since $\widetilde a_{ab}$ and $\widetilde c_a/24$ are
rational with bounded denominators).  For any $M^a$ in this kernel,
$R_a,R_0$ are then fixed by \eqref{eq:L2-PFV-integrality}.
Define
\begin{equation}\label{eq:L2-PFV-N-p}
    N_{ab}\coloneqq \widetilde\kappa_{abc}M^c,
    \qquad
    p^a\coloneqq N^{ab}K_b\, ,
\end{equation}
where $N^{ab}$ is the inverse of $N_{ab}$.  Via mirror symmetry the
rational vector $p=p^a\widetilde\omega_a$ is viewed as a direction in
the mirror K\"ahler cone.  The ordinary
derivatives of the polynomial superpotential are
\begin{equation}\label{eq:L2-PFV-dWpoly}
    \partial_{z^a}W_{\rm poly}=N_{ab}z^b-\tau K_a\, ,
    \qquad
    \partial_\tau W_{\rm poly}=-K_az^a\, .
\end{equation}
Thus $\partial_{z^a}W_{\rm poly}=0$ forces the complex line
\begin{equation}\label{eq:L2-PFV-line}
    z^a=p^a\tau\, .
\end{equation}
On this line one has
\begin{equation}
    W_{\rm poly}
      =-\frac12\tau^2K_ap^a\, ,
    \qquad
    \partial_\tau W_{\rm poly}
      =-\tau K_ap^a\, .
\end{equation}
Consequently the polynomial contribution satisfies
\begin{equation}
    W_{\rm poly}=\partial_\tau W_{\rm poly}
      =\partial_{z^a}W_{\rm poly}=0
\end{equation}
along $z^a=p^a\tau$, provided that
\begin{equation}\label{eq:L2-PFV-null-condition}
    K_ap^a=K_aN^{ab}K_b=0\, .
\end{equation}
Since $W_{\rm poly}=0$ on this locus, the covariant
$F$-terms $D_IW_{\rm poly}$ reduce to the ordinary derivatives
displayed above.

The PFV conditions are therefore the following arithmetic and
geometric requirements:
\begin{equation}\label{eq:L2-PFV-conditions}
    \det N\neq0,\qquad
    p=p^a\widetilde\omega_a\in\KC(\widetilde X),\qquad
    K_ap^a=0,\qquad
    \widetilde a_{ab}M^b\in\mathbb Z,\qquad
    M^a\widetilde c_a\in24\mathbb Z\, .
\end{equation}
The condition $p=p^a\widetilde\omega_a\in\KC(\widetilde X)$ is not cosmetic: it says
that the line $z^a=p^a\tau$ remains in the LCS region when
$\operatorname{Im}\tau$ is large.

\medskip
\noindent\emph{Diophantine viewpoint.}
For a mathematician, \eqref{eq:L2-PFV-conditions} is perhaps the
most interesting line in the construction.  The unknowns are not
continuous parameters but integral flux vectors $M^a,K_a$, together
with the induced rational vector $p^a=N^{ab}K_b$.  The equations
combine a quadratic null condition $K_aN^{ab}K_b=0$, congruence
conditions coming from the integral symplectic basis, the open cone
condition $p\in\KC(\widetilde X)$, and the D3-tadpole bound.
Thus the small-$W_0$ problem is, at this stage, a constrained
Diophantine search whose coefficients are the classical and
enumerative data of the Calabi--Yau.  The physics then asks for
solutions that additionally place the instanton expansion in a
controlled region.




After imposing
\eqref{eq:L2-PFV-conditions}, the only surviving contribution along
the flat direction is
\begin{equation}\label{eq:L2-PFV-Weff}
    W_{\rm eff}(\tau)
      =
      -\frac{1}{(2\pi)^2}
      \sum_{\widetilde{\mathbf q}\in\MC(\widetilde X)}
      n_{\widetilde{\mathbf q}}\,
      \widetilde q_aM^a\,
      \operatorname{Li}_2
      \!\left(
        \mathrm{e}^{2\pi\mathrm{i}\,\widetilde q_ap^a\tau}
      \right)\, .
\end{equation}
This is the promised bridge from the mirror-period computation to
flux vacua: once the truncated GV dictionary is known,
\eqref{eq:L2-PFV-Weff} is an explicit computable holomorphic function
of one variable.  The
racetrack discussion below is the elementary mechanism by which
several leading terms in this series can stabilise $\tau$ and
leave an exponentially small value of $W_0$.

\subsubsection*{The racetrack mechanism, made concrete}

The PFV mechanism produces a residual effective superpotential along
the flat direction $\tau$ (the variable already used in
\eqref{eq:L2-PFV-Weff}) of the form
\begin{equation}\label{eq:racetrack}
    W_{\rm eff}(\tau) \;=\; c\,\left (\mathrm{e}^{2\pi\mathrm{i}\,p_1 \tau}
        \,+\, A\,\mathrm{e}^{2\pi\mathrm{i}\,p_2 \tau}\right )
        \,+\, \cdots \,,
\end{equation}
where the rationals $p_1=\widetilde q^{(1)}_a p^a$ and
$p_2=\widetilde q^{(2)}_a p^a$ are built from the two leading curve
classes $\widetilde{\mathbf q}^{(1)},\widetilde{\mathbf q}^{(2)}$
and the rational vector $p^a=N^{ab}K_b$ of
\eqref{eq:L2-PFV-N-p}; they need not be integers.  The
prefactors $c, A$ are determined by the leading genus-zero
GV invariants of the mirror $\widetilde X$ that enter
\eqref{eq:L2-PFV-Wexp}.  Solving $\partial_{\tau}
W_{\rm eff}=0$ for $\tau$ and substituting back, the
residual value $|W_0|$ inherits the exponential smallness of the
worldsheet-instanton coefficients $c, A$ together with a fractional
power of their ratio controlled by $1/|p_2-p_1|$.
For example, if phases are aligned and
$p_2>p_1>0$ along the positive-imaginary saddle, the equation
$\partial_\tau W_{\rm eff}=0$ gives
\begin{equation}
    \mathrm{Im}\,\tau
    \simeq
    \frac{1}{2\pi(p_2-p_1)}
    \log\left|\frac{A p_2}{c p_1}\right|\, ,
\end{equation}
and substituting back gives schematically
\begin{equation}
    |W_0|
    \simeq
    |c|\left|1-\frac{p_1}{p_2}\right|
    \left|\frac{c p_1}{A p_2}\right|^{p_1/(p_2-p_1)}\, .
\end{equation}
The prefactors and phases depend on the chosen two leading terms,
but the combinatorial choice of flux quanta controls both the
exponent and the integer coefficients entering through the GV
invariants in $c, A$; see \cite{Demirtas:2019sip} for the exact
dependence.

\begin{example}[Two-modulus PFV example \cite{Demirtas:2019sip}]
                \label{ex:racetrack}
Consider the degree-$18$ hypersurface in
$\mathbb{P}^4_{[1,1,1,6,9]}$.  The full Calabi--Yau threefold has
$h^{2,1}=272$, but one restricts to the
$G=\mathbb Z_6\times\mathbb Z_{18}$-invariant locus of
\cite{Giryavets:2003vd}.  For $G$-invariant fluxes the F-term
equations are consistently restricted to this locus, and the relevant
periods are those of an effective two-modulus Greene--Plesser mirror
\cite{Greene:1990ud}.  This is the example used in
\cite{Demirtas:2019sip}; see also \cite{Candelas:1993dm} for the
two-parameter mirror-symmetry data.  The orientifold choice used there
has D3-charge budget $Q_{\mathrm{D3}}=138$.  We shall write the two
LCS coordinates as $z^1,z^2$.


In the notation of \eqref{eq:L2-F-poly-CMS}, the classical data are
\begin{equation}
    \widetilde\kappa_{111}=9,\qquad
    \widetilde\kappa_{112}=3,\qquad
    \widetilde\kappa_{122}=1,
\end{equation}
together with
\begin{equation}
    \widetilde a
    =
    \frac12
    \begin{pmatrix}
        9 & 3\\
        3 & 0
    \end{pmatrix},
    \qquad
    \widetilde c
    =
    \begin{pmatrix}
        102\\
        36
    \end{pmatrix}.
\end{equation}
The first instanton corrections to the prepotential are, in the
normalisation of \cite{McAllister:2025qwq} \S\,6.2,
\begin{equation}
    (2\pi\mathrm{i})^3F_{\rm inst}
    =
    F_1+F_2+\cdots ,
\end{equation}
where
\begin{align}
    F_1
    &=
    -540\,\mathrm{e}^{2\pi\mathrm{i}z^1}
    -3\,\mathrm{e}^{2\pi\mathrm{i}z^2},\\
    F_2
    &=
    -\frac{1215}{2}\,\mathrm{e}^{4\pi\mathrm{i}z^1}
    +1080\,\mathrm{e}^{2\pi\mathrm{i}(z^1+z^2)}
    +\frac{45}{8}\,\mathrm{e}^{4\pi\mathrm{i}z^2}.
\end{align}
The important numerical feature is already visible in $F_1$: the two
one-instanton coefficients are $540$ and $3$.  The racetrack will use
this hierarchy together with a flux choice that makes the two
exponents close.

Let $M=(M^1,M^2)^{\top}$ and $K=(K_1,K_2)^{\top}$ be the two flux
vectors in the PFV ansatz \eqref{eq:L2-PFV-flux-ansatz}.  With the
intersection numbers displayed above, the null condition
$K_aN^{ab}K_b=0$ becomes the explicit Diophantine relation
\begin{equation}
    M^1
    =
    \frac{M^2K_2(2K_1-3K_2)}
         {(K_1-3K_2)^2}\, .
\end{equation}
When this relation is obeyed and the integrality conditions
\eqref{eq:L2-PFV-integrality} hold, the flat PFV direction is
\begin{equation}
    \begin{pmatrix}
        z^1\\
        z^2
    \end{pmatrix}
    =
    \tau\,
    \frac{K_1-3K_2}{M^2}
    \begin{pmatrix}
        -K_2/K_1\\
        1
    \end{pmatrix}.
\end{equation}
Thus the problem has become arithmetic: one searches for integer
vectors $M,K$ satisfying this rational equation, the congruence
conditions needed for $R_0,R_a\in\mathbb Z$, the K\"ahler-cone
condition on $p$, and the tadpole bound.

Substituting the flat direction into the one-instanton terms gives a
two-term racetrack
\begin{equation}
    W_{\rm eff}(\tau)
    =
    c\left(
      \mathrm{e}^{2\pi\mathrm{i}p_1\tau}
      +A\,\mathrm{e}^{2\pi\mathrm{i}p_2\tau}
    \right)+\cdots ,
\end{equation}
where $p_1=p^1$ and $p_2=p^2$ for the two leading curve classes.
Solving $\partial_\tau W_{\rm eff}=0$ gives
\begin{equation}
    \mathrm{e}^{2\pi\mathrm{i}(p_1-p_2)\langle\tau\rangle}
    =
    -\,A\,\frac{p_2}{p_1},
    \qquad
    \langle\tau\rangle
    =
    \frac{1}{2\pi\mathrm{i}}\,
    \frac{1}{p_1-p_2}
    \log\!\left(-A\,\frac{p_2}{p_1}\right),
\end{equation}
and therefore
\begin{equation}
    W_{\rm eff}(\langle\tau\rangle)
    =
    c\,\frac{p_2-p_1}{p_2}
    \left(-A\,\frac{p_2}{p_1}\right)^{p_1/(p_1-p_2)}.
\end{equation}
The small number is produced when $|p_1-p_2|\ll p_2$: the two
instantons then compete at a large value of $\operatorname{Im}\tau$,
and the value of the superpotential is exponentially suppressed.

The explicit flux choice used in \cite{Demirtas:2019sip} is
\begin{equation}
    M=
    \begin{pmatrix}
        -16\\
        50
    \end{pmatrix},
    \qquad
    K=
    \begin{pmatrix}
        3\\
        -4
    \end{pmatrix}.
\end{equation}
The associated flat-direction vector is
\begin{equation}
    p=
    \begin{pmatrix}
        2/5\\
        3/10
    \end{pmatrix},
    \qquad
    z^1=\frac25\,\tau,\qquad
    z^2=\frac{3}{10}\,\tau .
\end{equation}
The flux contribution to the D3 tadpole is
\begin{equation}
    Q_{\rm flux}
    =
    -\frac12\,M^aK_a
    =
    124
    \leq 138.
\end{equation}
The two-instanton racetrack in \eqref{eq:racetrack} has
\begin{equation}
    A=-\frac{5}{288},
    \qquad
    c=
    -\sqrt{\frac{2}{\pi}}\,
    \frac{8640}{(2\pi\mathrm{i})^3}.
\end{equation}
Solving the full F-term equations gives
\begin{equation}
    \langle\tau\rangle=6.856\,\mathrm{i},\qquad
    \langle z^1\rangle=2.742\,\mathrm{i},\qquad
    \langle z^2\rangle=2.057\,\mathrm{i},
\end{equation}
and
\begin{equation}
    |W_0|=2.037\times10^{-8}.
\end{equation}
The accompanying GitHub repository contains a demo notebook,
\texttt{jaxvacua\_demo.ipynb} in the \texttt{notebooks/} directory,
which reconstructs this PFV input from the flux vectors above and
solves the corresponding F-term equations numerically.
The two-instanton terms displayed above are smaller by
$O(10^{-5})$ at this solution, so the two-instanton racetrack already
captures the scale of the answer.  


\end{example}

\paragraph{What this example demonstrates.}
This example shows how mirror pairs constructed from polytope
triangulations can turn finite combinatorial input into integer
enumerative data, and how those integers can then produce very small
numbers in string-theoretic quantities.  The small value of
$W_0$ is not obtained by tuning continuous parameters: it comes from
GV invariants in the instanton expansion, combined with integer flux
quanta through the PFV Diophantine conditions.  This is highly
non-trivial, and it is one of the key inputs for the next lecture,
where the goal is to compute the full vacuum energy; the same
suppression mechanism is what makes parametrically tiny scales
available in the effective theory.



\paragraph{Lecture~2 punchlines.}
The second lecture has one computational message.  To use
fluxes in a Calabi--Yau compactification one must compute the periods
of the holomorphic three-form, because
\begin{equation}
    W_{\rm flux}=\int_X\Omega\wedge G_3
\end{equation}
is a flux-weighted period integral.  The integers specifying
$F_3$ and $H_3$ choose the linear combination; mirror symmetry and
toric geometry compute the functions that appear in that linear
combination.

The mirror computation translates the same period vector
into the large volume $A$-model prepotential on $\widetilde X$.  Its
instanton part is organised by genus-zero GV invariants,
so these integers enter the four-dimensional superpotential as
actual coefficients, not merely as auxiliary enumerative data.  The
pipeline established in the lecture is therefore
\begin{equation}
\begin{aligned}
    \text{toric hypersurface data}
    &\quad\Longrightarrow\quad
    \text{periods and mirror map}\\
    &\quad\Longrightarrow\quad
    \text{GV invariants}
    \quad\Longrightarrow\quad
    W_{\rm flux}\, .
\end{aligned}
\end{equation}
The PFV computation at the end of the lecture is the first
application of this pipeline.  It chooses fluxes so that the
polynomial LCS contribution to $W_{\rm flux}$ has a flat solution
locus, and then uses the GV-controlled instanton terms to lift that
locus and produce a small value of $W_0$.  This is precisely the
calculation illustrated in the companion notebook
\texttt{notebooks/jaxvacua\_demo.ipynb}.  Lecture~3 starts from this
flux-sector output and asks what is still required to obtain an
actual vacuum energy.
